\subsection{BARYCENTRIC CALCULUS}

PROPOSITION 1. Let \((x_{t})_{t\in I}\) be a family of points in an affine space \(\mathbf{E}\) and \((\lambda_t)_{t\in I}\) a family of elements of \(\mathbf{K}\) of finite support such that \(\sum_{i\in I}\lambda_i = 1\) (resp. \(\sum_{i\in I}\lambda_i = 0\) ). If \(a\) is any point of \(\mathbf{E}\) , the point \(x\in \mathbf{E}\) defined by

\[
x - a = \sum_ {i \in I} \lambda_ {i} (x _ {i} - a)
\]

(resp. the free vector \(\sum_{i\in I}\lambda_i(x_i - a))\) is independent of the point considered.

If \(a'\) is another point of \(E\) , then

\[
\sum_ {i} \lambda_ {i} (x _ {i} - a ^ {\prime}) = \sum_ {i} \lambda_ {i} ((x _ {i} - a) + (a - a ^ {\prime})) = \sum_ {i} \lambda_ {i} (x - a) + \left(\sum_ {i} \lambda_ {i}\right) (a - a ^ {\prime}).
\]

If \(\sum_{i}\lambda_{i} = 1\) , then \(\sum_{i}\lambda_{i}(x_{i} - a) = (x - a) + (a - a') = x - a'\) ; if \(\sum_{i}\lambda_{i} = 0\) , then \(\sum_{i}\lambda_{i}(x - a') = \sum_{i}\lambda_{i}(x - a)\) ; whence the proposition.

Under the conditions of Proposition 1, the point x defined by

\[
x - a = \sum_ {i \in I} \lambda_ {i} (x _ {i} - a)
\]

\(\left(\text{resp. the free vector} \sum_{t \in I} \lambda_t(x_t - a)\right)\) is denoted by \(\sum_{t \in I} \lambda_t x_t\) .

Thus in particular the notation \(b - a\) introduced in no. 1 is recovered. When \(\sum_{i} \lambda_{i} = 1\) , the point \(x = \sum_{i} \lambda_{i} x_{i}\) is called the barycentre of the points \(x_{i}\) given the masses \(\lambda_{i}\) .

Given \(m\) points \(a_1, \ldots, a_m\) of \(\mathbf{E}\) , whose number \(m\) is not a multiple of the characteristic of \(\mathbf{K}\) (V, § 1), the point \(g = \sum_{i=1}^{m} \frac{1}{m} a_i\) is called (by an abuse of language) the barycentre of the points \(a_i\) ( \(1 \leqslant i \leqslant m\) ) (for \(m = 2\) , we say ``midpoint'' instead of ``barycentre''); it is characterized by the relation

\[
\sum_ {i = 1} ^ {m} \left(a _ {i} - g\right) = 0.
\]

